% META: {"id": "TNSI-BDD-EVAL-B", "chapitre": "TNSI-BASES-DE-DONNEES", "type_objet": "evaluation", "capacites": ["T-BDD-02", "T-BDD-03D", "T-BDD-03E", "T-BDD-03G", "T-BDD-03H"], "mode_creation": "ex_nihilo", "sources_inspiration": [], "duree_min": 45, "status": "generated"}
\section*{Évaluation B --- Bases de données}
\textit{Durée : 45 minutes. Ordinateur avec un SGBD (SQLite ou équivalent) autorisé.}

\baremeIndicatif{Ex. 1 : 6 pts (SGBD) ; Ex. 2 : 14 pts (JOIN, ORDER BY, UPDATE, DELETE)}

\bigskip

\textbf{Exercice 1} --- \textit{Services d'un SGBD} \hfill (6 points)

\begin{enumerate}
  \item (3 pts) Citer les quatre services rendus par un SGBD relationnel.
  \item (3 pts) Expliquer pourquoi un simple fichier CSV ne peut pas remplacer un SGBD pour gérer les réservations d'une salle de sport ouverte à plusieurs guichets simultanément.
\end{enumerate}

\bigskip

\textbf{Exercice 2} --- \textit{JOIN, ORDER BY, UPDATE, DELETE} \hfill (14 points)

On donne :
\begin{sql}
CREATE TABLE Employe(id INTEGER PRIMARY KEY, nom TEXT, id_service INTEGER);
CREATE TABLE Service(id INTEGER PRIMARY KEY, nom TEXT);
INSERT INTO Service VALUES (1, 'Comptabilite'), (2, 'Informatique');
INSERT INTO Employe VALUES (1, 'Ben Ali', 2), (2, 'Chraibi', 1), (3, 'Haddad', 2);
\end{sql}
\begin{enumerate}
  \item (5 pts) Écrire la requête affichant le nom de chaque employé avec le nom de son service, triée par nom d'employé croissant.
  \item (4 pts) Écrire la requête qui change le service de l'employé \texttt{'Ben Ali'} vers le service \texttt{1}.
  \item (5 pts) Écrire la requête qui supprime tous les employés du service \texttt{1} (après la modification de la question précédente). Que se passerait-il si l'on omettait la clause \texttt{WHERE} ?
\end{enumerate}

% BEGIN-VERIFY
% import sqlite3
% conn = sqlite3.connect(":memory:")
% conn.execute("CREATE TABLE Employe(id INTEGER PRIMARY KEY, nom TEXT, id_service INTEGER)")
% conn.execute("CREATE TABLE Service(id INTEGER PRIMARY KEY, nom TEXT)")
% conn.executemany("INSERT INTO Service VALUES (?,?)", [(1, "Comptabilite"), (2, "Informatique")])
% conn.executemany("INSERT INTO Employe VALUES (?,?,?)", [(1, "Ben Ali", 2), (2, "Chraibi", 1), (3, "Haddad", 2)])
% conn.commit()
% cur = conn.execute("""SELECT Employe.nom, Service.nom FROM Employe
%     JOIN Service ON Employe.id_service = Service.id ORDER BY Employe.nom ASC""")
% assert cur.fetchall() == [("Ben Ali", "Informatique"), ("Chraibi", "Comptabilite"), ("Haddad", "Informatique")]
% conn.execute("UPDATE Employe SET id_service = 1 WHERE nom = 'Ben Ali'")
% conn.commit()
% assert conn.execute("SELECT id_service FROM Employe WHERE nom = 'Ben Ali'").fetchone() == (1,)
% conn.execute("DELETE FROM Employe WHERE id_service = 1")
% conn.commit()
% assert conn.execute("SELECT COUNT(*) FROM Employe").fetchone()[0] == 1
% cur = conn.execute("SELECT nom FROM Employe")
% assert cur.fetchall() == [("Haddad",)]
% END-VERIFY
